\section{Scenario, Candidate Actions, and Configuration}
\label{app:configuration}

\subsection{Population and Target Selection}

The input scenario contains 84,110 people. Sampling sorts person identifiers, shuffles them with Java Random using seed 6021, and retains a fixed 20\% subset with all its plans, yielding 16,822 simulated travelers. Network flow and storage capacity factors are 0.2 and $0.2^{0.75}\approx0.2991$, respectively \cite{mersiniScalingMATSimTraffic2026,ben-dorPopulationDownscalingMultiagent2021}. Population sampling, synthetic profiles, information groups, and controlled-roster selection use seed 6021; the main simulation seed is 268.

Targets are selected from original plans using destination-activity start times in $[0,32400)$ seconds, corresponding to starts before 09:00. Selection covers all activity purposes and requires one car leg per target trip; it imposes no departure-window filter. Missing or invalid planned arrival times are excluded rather than imputed. The selected cohort contains 11,793 trips by 11,793 people across 2,877 OD pairs. Membership remains fixed throughout the experiment. The remaining 5,029 people and non-target trips of controlled people retain MATSim decision mechanisms. Work-arrival evaluation uses the fixed subset of 11,332 targets with work deadlines. The 09:00 selection threshold does not substitute for each trip's deadline.

\subsection{Fixed Candidate Sets}

Runs reuse the same precomputed BFS-LE candidate file \cite{rieser2013route}, validating topology, costs, overlap, and the shortest-route candidate without regenerating or reordering alternatives. The archived settings specify up to four candidates per OD pair and maximum pairwise overlap 0.85. For free-flow time, an alternative passes the detour check if its ratio to the reference is at most 2.5 or its excess time is at most 300 seconds. The analogous distance limits are 2.5 and 2,000 meters; both time and distance checks must pass. The exact overlap implementation and original candidate-file hash belong to the reproducibility archive.

Departure candidates span $-5400$ to $+1800$ seconds relative to the original planned departure, at 900-second intervals. Options are clipped against the original origin-activity start, day zero, and the exclusive simulation end; the original departure is retained. The arrival deadline is not a hard feasibility filter. Subsequent choices do not shift the original time anchor. Of the targets, 2,780 have one route and 9,013 have multiple routes; the former cannot contribute route switches within the fixed candidate space.

\begin{table}[htbp]
\centering\small
\caption{Principal archived settings for the 100-day experiment. Strategy weights are configuration weights, not probabilities inferred from their sum.}
\label{tab:configuration}
\renewcommand{\arraystretch}{1.15}
\begin{tabularx}{\linewidth}{@{}>{\raggedright\arraybackslash}p{0.35\linewidth}X@{}}
\toprule
Setting & Value \\
\midrule
MATSim version & 2026.0 \\
Population / simulation seeds & 6021 / 268 \\
Initial strategy weights & \texttt{ChangeExpBeta}: 0.7; \texttt{ReRoute}: 0.1; \texttt{TimeAllocationMutator}: 0.1 \\
Innovation cutoff & Log Day90, equivalent to paper day 91; existing-plan selection continues \\
Simulation interval for charging & $[0,108000)$ seconds \\
Manager opportunities & Paper days 21, 51, and 81 \\
Toll action limits & At most 40 distinct car links; finite nonnegative charge per entry; no preset price cap or increment \\
Recent personal observations & Up to three preceding days \\
Personal summary limit & 500 characters \\
Manager summary limit & 1,000 English words; target length 400--600 words, with no enforced minimum \\
Information groups & 5,897 informed and 5,896 uninformed; fixed person-level assignment \\
\bottomrule
\end{tabularx}
\end{table}

The effective innovation cutoff is established from the run logs and final configuration, rather than inferred solely from the initial strategy entries. Both background replanning and MATSim-agent target decisions are affected; LLM and MNL target choices continue under their respective decision mechanisms.

\section{Information, Memory, and Execution Protocol}
\label{app:implementation}

\subsection{Information Access and Initialization}

Announcement visibility is assigned at person level: all trips of a controlled person share the visibility assignment. The LLM decides only the designated target trip; other trips retain background routing and time-adjustment mechanisms. For uninformed people, current toll information remains unknown, while background routing can retain personal link/time-bin observations. Thus, sharing an information assignment does not mean that every trip is decided by an LLM. People outside the treatment use the background information mechanism.

Informed target decisions receive current candidate-related toll estimates. Uninformed decisions receive null current costs, not zero costs. Both groups receive static candidate information, personal recent outcomes, and their own summaries, without network-wide realized traffic or unchosen-route outcomes. LLM Traveler explicitly communicates that charging is disabled. In the combined configuration, informed and uninformed inputs remain distinct even when the active policy is empty. Initial factual histories and self-written summaries are empty; synthetic profiles represent preference assumptions rather than past experiences.

Managers receive public network information, aggregate reference-plan OD diagnostics, completion statistics, aggregate link traffic, and their own policy and memory histories. They do not receive private traveler memories, profiles, or individual payments. The archived manager projection retains the latest two executed policy-version effects and pools older days; it includes aggregate recent/previous statistics over windows of up to seven days. All-car links remain represented in the public network, while detailed traffic projection prioritizes high-delay or failed links and currently or recently charged links. This projected manager history is distinct from the travelers' three-day personal window.

\subsection{Memory Updates and Output Validation}

Traveler summaries are fully replaced using prior memory and newly available feedback. They remain private to the person and target trip and are not shared across agents. No separate exact-payment history is available to the LLM outside its bounded observations and summary. The archive retains full records for research, but does not serve as an unrestricted retrieval memory. Executable route choices are checked against candidates; references to unavailable route identifiers within summary text are logged diagnostically rather than automatically rewritten. Semantic review was disabled, so syntactic validity does not certify factual correctness.

Traveler output validation checks exact JSON, permitted route and departure values, and character limits. There are at most three output attempts, comprising the initial attempt and up to two retries. Truncated responses are rejected without salvaging partial JSON or inventing actions. Persistent traveler-decision or memory failure terminates the experiment rather than silently reusing a stale summary. Late arrival and noncompletion are execution outcomes, not automatically output-format errors.

Manager memory is updated after each completed day, including toll-free days, independently of scheduled policy decisions. A memory update cannot change the policy. Failed or invalid management-memory calls retain the previous memory, and the next successful update includes the missing days. Policy episodes are identified by effective prices, so an unchanged policy version does not create a new price intervention. Detailed transport-error and policy-fallback handling must be read together with the archived implementation; output repair counts alone do not characterize those paths.

\subsection{Plan Execution and Billing}

Scheduled management decisions precede traffic execution. MATSim replanning is followed by writing validated controlled route and departure choices into target plans. Other selected-plan activity ends are retained. Original, selected, and realized departure times are stored separately, since earlier activities, queues, and previous trips can delay actual execution.

Tolls are charged on eligible car-link entries, including the origin partial link under the archived billing convention. Repeated traversals incur repeated charges, while duplicate notifications of the same entry are ignored. Charges apply on the half-open interval $[0,108000)$ seconds and remain constant within a day. Payments already incurred are retained if a trip subsequently fails to complete. Candidate cost estimates and actual payments are separate quantities. A valid policy replaces the entire rule set; an empty policy removes charges. Link-delay statistics use completed full-link traversals and therefore have different event boundaries from toll collection and trip-duration measurement.

\section{Decision Models and Prompt Specification}
\label{app:decision-models}

\subsection{Constrained MNL Baseline}

The archived baseline version is \texttt{morning-joint-mnl-v1}. For candidate action $a=(r,\tau)$, its deterministic utility is
\begin{equation}
 V_{i,d}(a)=-0.3\,\widehat{T}_{i,d}(a)-0.3\,\widehat{C}_{i,d}(a)
           -0.9\,\widehat{L}_{i,d}(a)-0.05\,|\tau-\tau_i^0|,
 \label{eq:appendix-mnl}
\end{equation}
where time, predicted lateness, and departure deviation are measured in minutes and money in scenario units. $\tau_i^0$ is the original departure. Choice probabilities are
\begin{equation}
 P_{i,d}(a)=\frac{\exp V_{i,d}(a)}{\sum_{b\in\mathcal{R}_i\times\mathcal{T}_i}\exp V_{i,d}(b)}.
\end{equation}
The softmax scale is one; no additional logit multiplier is applied. Coefficients are homogeneous engineering assumptions, with no numerical mapping from qualitative profiles and no calibration to observed behavior.

When matching completed personal observations exist within the three-day window, expected travel time is an equal blend of static free-flow time and the mean observed car time for the same route and selected-departure bin. Otherwise it equals free-flow time. Bins are 900 seconds wide and use selected, not delayed actual, departure times. Failed trips are excluded from duration and payment means but remain in audit records. The supplied parameter file sets unknown-toll mode to \texttt{reject}; its illustrative 0.1-unit per-link prior is inactive. The reported MNL experiment has no tolls and does not evaluate that prior or an uninformed pricing response.

Choice seeds are derived from a SHA-256 hash of the baseline version, seed, simulation day, and trip identifier, using the first 64 bits for a local Java Random instance and a canonical departure/route ordering. Candidate costs, utilities, probabilities, and seeds are archived. This baseline uses neither natural-language memory summaries nor network-wide traffic observations.

\subsection{LLM Generation and Input Templates}

Table~\ref{tab:llm-settings} reports settings recorded for the combined run. These are configured model identifiers, not a claim that remote model weights or service implementations are immutable. Generation temperature and top-$p$ are not specified in the inspected settings; their effective service defaults should not be inferred to be zero or identical across runs.

\begin{table}[htbp]
\centering\small
\caption{Archived generation settings in LLM Traveler + Manager. Token limits are per-response maxima, not observed consumption. All listed request timeouts are 300 seconds.}
\label{tab:llm-settings}
\begin{tabularx}{\linewidth}{@{}>{\raggedright\arraybackslash}X >{\raggedright\arraybackslash}p{0.26\linewidth}r >{\raggedright\arraybackslash}p{0.19\linewidth}@{}}
\toprule
Task & Model ID & Max. tokens & Response format \\
\midrule
Traveler choice & \texttt{qwen3.5-9b} & 1,024 & JSON schema \\
Personal memory & \texttt{qwen3.5-9b} & 1,024 & JSON schema \\
Manager policy & \texttt{deepseek-flash} & 8,192 & JSON object \\
Manager memory & \texttt{deepseek-flash} & 2,048 & JSON object \\
\bottomrule
\end{tabularx}
\end{table}

Thinking is disabled in the recorded generation settings. The traveler client specifies concurrency 256, a 120-character reason limit, and a 900-second stall threshold. These client settings do not determine wall-clock performance on different deployments.

The traveler and personal-memory template identifiers are \texttt{morning-information-v9-continuous-memory} and \texttt{morning-reflection-v9-continuous-memory}. Manager policy and memory use \texttt{morning-pricing-v10-daily-memory} and \texttt{morning-manager-memory-v1-daily}. Inputs use JSON-array tables with declared columns, exact numeric seconds, and explicit nulls. Common rules precede fixed personal facts and daily observations; manager topology precedes policy and traffic information. The manager policy requires \texttt{rules} and \texttt{decision}, while its independent reflection requires \texttt{memory}.

% AUTHOR INSERT: Verbatim system/user templates, output JSON schemas, one
% redacted instantiated example per role, and template hashes from the exact
% run artifacts. The summaries above are not substitutes for original prompts.
% AUTHOR VERIFY: Effective temperature/top-p and served-model revision for each
% deployment; precise work-deadline extraction and overlap implementation.

\section{Measurement and Aggregation}
\label{app:measurement}

All five runs contribute the same paper days 1--100 (log Days0--99). Daily outcome summaries are averaged over this window. Relative changes are computed as the ratio of window means minus one; they do not average daily percentages. Lateness differences are percentage points. Switching rates use 99 valid adjacent-day transitions. For cutoff diagnostics, log Day90 is treated separately; the post-cutoff comparison uses Days91--99.

Total daily car time sums car-departure-to-car-arrival durations for the entire population, including waiting before network entry. A complete total is undefined if a planned car trip is incomplete. Completed-target travel-time distributions exclude incomplete trips; completion and work-arrival denominators retain the original cohort. For work targets, on-time completion, late completion, and noncompletion exhaust the original set. All planned car trips completed in the reported experiment, so this missing-total rule did not remove any reported day.

Route identity uses the complete ordered link sequence, not a route label that may remain unchanged under MATSim replanning. Candidate labels such as \texttt{c0} are local to an OD pair and cannot identify a network-wide route. Net OD-route turnover divides the sum of absolute changes in route counts by $2N$, distinguishing aggregate redistribution from individual switching. Multi-route switching uses only the 9,013 eligible candidate-based targets; this denominator is not applied to MATSim agents.

Mean daily P95 averages each day's target-trip 95th percentile. Individual savings compare the same target person's 100-day mean times. Early arrival averages $\max(\mathrm{deadline}-\mathrm{arrival},0)$ over all work targets, including zeros. Selected departure offsets and actual departures remain separate. No monetary value is assigned to early arrival or departure inconvenience.

Link-delay comparisons aggregate original link records to full days before combining runs with different reporting intervals. Completed-link delay excludes pre-network waiting. Target/other-trip decompositions and waiting/post-entry decompositions each recover total car time, but must not be added to one another. Target payments and full-day personal payments likewise have different scopes. Daily smoothing uses a trailing seven-day mean, not a confidence interval. No independent-sample significance claims are made from successive days or travelers within the same run.

\section{Execution Audit and Computational Workload}
\label{app:audit-cost}

Output repair was required for 2,433 finalized travel-decision records (0.206\%) and 2,943 personal-memory records (0.250\%). All personal memory chains were continuous, submitted memories remained within 500 characters, and no finalized travel decision was blank. The three management decisions were valid; one invalid management-memory JSON response was repaired through retry, with all 100 daily updates completed. These checks establish consistency of accepted outputs and state updates, without implying that every initial generation was valid or that generated reasons and memories were semantically correct.

\begin{table}[htbp]
\centering
\small
\caption{Recorded LLM inference workload over complete 100-day runs. Counts include retries and repairs across roles; tokens are service-reported totals. Zero LLM calls does not imply zero simulation cost.}
\label{tab:cost}
\begin{tabularx}{\linewidth}{@{}Xrr@{}}
\toprule
Condition & API calls & Tokens \\
\midrule
MATSim & 0 & 0 \\
MNL & 0 & 0 \\
LLM Traveler & 2,362,964 & 5,713,538,990 \\
MATSim + Manager & 109 & 872,774 \\
LLM Traveler + Manager & 2,364,090 & 5,451,720,152 \\
\bottomrule
\end{tabularx}
\end{table}

The combined run contains 1,181,736 traveler-decision calls, 1,182,250 personal-memory calls, three management-decision calls, and 101 management-memory calls. There are 1,179,300 finalized records for each personal task. The excess calls include repair and retry overhead and cannot be interpreted as hallucination counts. Logged tokens do not measure monetary cost, GPU energy, or end-to-end runtime. Semantic correctness of all memories and faithfulness of generated reasons have not been systematically established.

\section{Supplementary Behavioral and Management Results}
\label{app:supplementary}

\begin{table}[htbp]
\centering\small
\caption{Mean daily choice-change rates over 99 adjacent-day transitions. Overall rates use 11,793 targets; multi-route rates use 9,013 eligible targets. A dash denotes a candidate-space comparison that is not applicable. All entries are percentages.}
\label{tab:supp-switching}
\begin{tabularx}{\linewidth}{@{}Xrrr@{}}
\toprule
Condition & \shortstack{Overall\\route changes} & \shortstack{Multi-route\\changes} & \shortstack{Departure\\changes} \\
\midrule
MATSim & 5.26 & --- & 20.05 \\
MNL & 48.55 & 63.52 & 85.24 \\
LLM Traveler & 3.08 & 4.03 & 29.45 \\
MATSim + Manager & 5.43 & --- & 20.28 \\
LLM Traveler + Manager & 2.19 & 2.86 & 28.04 \\
\bottomrule
\end{tabularx}
\end{table}

Among multi-route targets, 48.98\% in LLM Traveler and 52.62\% in the combined configuration never switched during the 99 transitions. These proportions describe route persistence; single-route people are excluded and persistence is not treated as proof of learned habits.

The mean daily maximum 15-minute share of actual target departures was 20.71\% for LLM Traveler and 18.54\% for the combined configuration. Their mean shares of all-car link delay on the ten highest-delay links were 36.83\% and 38.92\%, respectively, ranking links independently each day. Thus, a less concentrated target departure distribution coexisted with more concentrated link delay and higher total car time in the combined configuration. These descriptive measures do not identify a social coordination mechanism.

Under seed 9261, MATSim + Manager increased total daily car time by 0.26\% relative to MATSim, while target time decreased by 0.18\% and work lateness by 0.06 percentage points. Under seed 268, the corresponding changes were $+1.86\%$, $+1.06\%$, and $+0.45$ percentage points. This supplementary comparison does not replicate the LLM traveler or information-group experiments.

\subsection{Archive and Reproduction Scope}

The analysis preserves run manifests, final effective configurations, input and executable hashes, candidate files, accepted actions, attempt logs, experience records, policy histories, and memory updates. Input matching does not imply executable equivalence: MATSim and MATSim + Manager share a build, while other builds differ. Final-output checks establish constraint and record consistency; the 4,000-context visibility audit covers supplied samples rather than every decision. Software artifacts, original prompts, and service-generation settings are needed in addition to the aggregate tables to reproduce the complete experiment.

% AUTHOR INSERT before release: repository URL and immutable revision, per-run
% executable hashes and model deployment identifiers, exact original prompt
% files, final configs, and instructions linking raw records to analysis outputs.
% \FloatBarrier
